%% ========== METHODS ==========
\section*{Methods}

\paragraph{Gaussian splat model.}
Each splat $k$ is parameterized by a center $\boldsymbol{\mu}_k \in \mathbb{R}^D$, a positive amplitude $a_k$ (via softplus activation), and a lower-triangular Cholesky factor $\mathbf{L}_k \in \mathbb{R}^{D \times D}$.
An optional sharpness exponent $s_k$ (via softplus, clamped to $[1.5, 4.0]$) generalizes the Gaussian shape; all results in this paper use the standard Gaussian ($s_k = 2$).
The model prediction at voxel $\mathbf{x}$ is
$\hat{V}(\mathbf{x}) = \sum_{k} a_k \exp\!\left(-\frac{1}{2}\|\mathbf{L}_k^{-1}(\mathbf{x} - \boldsymbol{\mu}_k)\|^{2}\right)$.
The Cholesky parameterization stores $D(D{+}1)/2$ elements per splat (6 for 3D) and guarantees positive-definite covariance by construction.
Amplitudes are clamped to $[0, \infty)$ via softplus with a configurable $\beta$ parameter.

\paragraph{Optimization.}
We use a per-splat Adam optimizer\cite{kingma2015adam} ($\beta_1{=}0.9$, $\beta_2{=}0.999$, $\epsilon{=}10^{-8}$) that maintains individual momentum and variance accumulators for each splat's parameters.
This enables dynamic operations (relocation, pruning, seeding of new splats) without disrupting the optimization state of existing splats.
The base learning rate is $0.01$ with dimensional compensation: for $D$-dimensional data, the effective learning rate is scaled by $D^{0.8} \times P_D / P_2$, where $P_D$ is the number of parameters per splat in $D$ dimensions.

The loss function is:
\begin{equation}
\mathcal{L} = \mathcal{L}_\text{recon} + \lambda_a \|\mathbf{a}\|_1 + \lambda_d \|\text{diag}(\mathbf{L})\|_1 + \lambda_b \mathcal{L}_\text{boundary}
\end{equation}
where $\mathcal{L}_\text{recon}$ is an asymmetric L1 loss that penalizes over-prediction $10\times$ more than under-prediction, $\lambda_a = 10^{-5}$ and $\lambda_d = 10^{-4}$ regularize splat amplitude and extent, and $\mathcal{L}_\text{boundary}$ penalizes splats whose centers fall outside the volume.

\paragraph{Dynamic operations.}
Every 500 iterations, the optimizer performs dynamic topology adjustment: (1) splats with amplitude below a threshold ($10^{-4}$ of the volume maximum) are pruned; (2) residual-guided relocation moves the lowest-amplitude splats to positions of peak reconstruction error; (3) Z-order (Morton code) sorting is applied to improve GPU cache locality.
The relocation step detects under-represented regions by computing $|\hat{V} - V|$, finding local maxima in the residual, and teleporting selected splats to these positions while preserving their momentum states.

\paragraph{Seed generation.}
Initial splat positions are determined by edge-based seeding: Sobel gradients are computed across all dimensions, local peaks in gradient magnitude are detected, and positions are deduplicated with a minimum separation criterion.
The number of seeds requested typically exceeds the final count, as pruning removes low-amplitude splats during optimization.
Seeding supports GPU acceleration (CUDA, MPS) for 10--50$\times$ speedup on large volumes.

\paragraph{Tiled fitting.}
For volumes exceeding GPU memory, \luxar partitions the volume into overlapping tiles on a regular grid with stride $T - L$, where $T$ is the tile size and $L$ is the overlap.
Each tile is fitted independently with the same optimizer configuration.
Splat amplitudes are modulated by a cosine (Hann) window $w(k) = \frac{1}{2}(1 - \cos(\pi k / (L+1)))$ applied along each overlap face, producing a partition-of-unity: adjacent windows sum to 1.0.
Boundary faces retain unit weight.
The constraint $L \leq T/2$ prevents triple overlap.
After fitting, tiles are concatenated into a single splat dataset.

\paragraph{Data encoding and spatial indexing.}
Fitted splats are written to Zarr\cite{miles2023zarr} archives with semantic encoding (7 data types, 4 compression modes) and Blosc compression (zstd level~3).
Spatial indexing uses compound ordering: discrete non-spatial dimensions (time, channel) form the most-significant bits, followed by a Morton (Z-order) curve over spatial coordinates.
This ordering ensures that spatially proximate splats are stored in adjacent chunks, enabling efficient view-frustum culling at load time.
Chunk-level axis-aligned bounding boxes are precomputed and stored as metadata.

\paragraph{Web viewer architecture.}
The \luxar viewer is a TypeScript application rendering via WebGL with Three.js.
Splat data is loaded progressively via HTTP range requests into a three-level cache: L2 (persistent OPFS, 2~GB), L1 (compressed RAM, 100~MB), L0 (decompressed GPU buffers, 200~MB).
Decompression uses WebAssembly (WASM) with support for up to 16 dimensions.
Rendering operates in 16-bit floating-point for HDR fidelity, with configurable tone mapping (ACES filmic, Reinhard, neutral, linear, Cineon, AgX), exposure, bloom (8-level mipmap), and anti-aliasing (FXAA, MSAA, SMAA, SSAA).
Adaptive device-pixel-ratio scaling maintains interactive frame rates.

For $n$-dimensional datasets, the viewer displays three spatial dimensions and navigates non-displayed dimensions via sliders (continuous) or toggles (categorical/boolean).
Gaussian splats are treated as hyperspheres: the effective display radius is $r_\text{eff} = \sqrt{r^2 - d^2}$, where $d$ is the Mahalanobis distance from the splat center to the current hyperplane.
Per-dimension affine transforms (scale, offset) and categorical permutations are supported via an inverse-query approach that transforms the query in O(1) rather than transforming all point positions.

\paragraph{Rate-distortion analysis.}
For each of the 12 datasets, we fitted Gaussian splat models at 10 splat counts logarithmically spaced from 1K to 512K.
All fits used conservative culling (99.9\% amplitude retention), 20,000 maximum iterations, early stopping patience of 500 iterations, L1 loss, and enabled dynamic operations.
PSNR was computed as $10 \log_{10}(1 / \text{MSE})$ on volumes normalized to $[0, 1]$.
SSIM\cite{wang2004image} was computed with a window size of 7 and Gaussian weighting.

\paragraph{Noise2Self analysis.}
We adapted the Noise2Self\cite{batson2019noise2self} blind-spot framework for Gaussian splat model selection.
For each dataset and splat count, 5\% of voxels were randomly selected and replaced with the median of their 3$\times$3$\times$3 donut neighborhood (excluding the center voxel).
Gaussian splats were fitted to this modified volume.
The train PSNR was computed on unmasked voxels against the original volume; the held-out PSNR was computed on the masked voxels against the original (pre-replacement) values.
The splat count at which the held-out PSNR peaks identifies the point beyond which additional splats memorize noise rather than capture signal.

\paragraph{Compression baseline comparison.}
We compared Gaussian splats against H.265 video compression (libx265, treating Z-stacks as video sequences) using the Noise2Self blind-spot framework\cite{batson2019noise2self} applied identically to both methods.
For each dataset, 5\% of voxels were randomly selected and replaced with the median of their 3$\times$3$\times$3 donut neighborhood (26-connected, excluding the center voxel).
Both GSplats and H.265 were then applied to this modified volume: GSplats were fitted to it (the optimizer sees donut-filled values at masked positions); H.265 encoded it at CRF values 0--51.
In both cases, the held-out PSNR was computed by comparing the output at masked positions against the \emph{original} (pre-replacement) values, measuring each method's ability to recover the true signal from surrounding context.
Both methods are spatially aware (GSplats: global mixture model; H.265: block-based prediction and DCT), making the blind-spot framework applicable to both.

\paragraph{Noise floor estimation.}
For each dataset, we estimated the noise standard deviation using three independent methods: (1) Laplacian of Gaussian MAD (median absolute deviation of Laplacian-filtered volume, scaled by $1/(\sigma_\text{MAD} \times 0.6745)$); (2) Haar wavelet MAD (finest-scale diagonal detail coefficients); (3) background percentile MAD (lowest 5\% intensity voxels).
The ensemble estimate is the median of the Laplacian and Haar estimates.
The noise floor PSNR is $10\log_{10}(1/\sigma^2)$.

\paragraph{Datasets.}
All datasets are from public repositories.
OpenCell spinning-disk confocal\cite{cho2022opencell}: MAP4-GFP (microtubules) and Hoechst (nuclei), 51$\times$600$\times$600 voxels; LMNB1, 112$\times$600$\times$600 voxels.
Mouse kidney confocal: DAPI (nuclei) and Phalloidin (actin), 16$\times$512$\times$512 voxels, acquired by G.~Buckley (Monash Micro Imaging), distributed via scikit-image\cite{vanderWalt2014scikit}.
Cells3D (hiPSC, Allen Institute\cite{viana2023allen}): nuclei and membrane channels, 60$\times$256$\times$256, via scikit-image\cite{vanderWalt2014scikit}.
3D cell culture (mouse ESC acini)\cite{blin2019nessys}: 236$\times$275$\times$271 voxels, from IDR (idr0062).
\emph{C.~elegans} embryo confocal\cite{murray2008celegans}: timepoint 100, 41$\times$512$\times$512, via the Cell Tracking Challenge\cite{maska2023celltracking}.
Tribolium embryo light-sheet: 991$\times$104$\times$965 voxels (cropped, multi-view deconvolved)\cite{preibisch2014efficient}.
All volumes were normalized to $[0, 1]$ intensity range before fitting.
